% Part: first-order-logic
% Chapter: completeness
% Section: maximally-consistent-sets

% महत्तम सुसंगत संचांची व्याख्या. संपूर्णतेसाठी आवश्यक असलेले त्यांचे गुणधर्म वापरलेल्या सिद्धता-प्रणालीच्या प्रकरणातील maximally-consistent-sets.tex मध्ये सिद्ध केले आहेत.

\documentclass[../../../include/open-logic-section]{subfiles}

\begin{document}

\olfileid{fol}{com}{mcs}
\olsection{\usetoken{P}{sentence} यांचे महत्तम सुसंगत संच}

\begin{defn}[महत्तम सुसंगत संच]
\ollabel{mcs}
!!{sentence}s चा संच~$\Gamma$ \emph{महत्तम सुसंगत} तेव्हा आणि केवळ तेव्हाच असतो, जेव्हा
\begin{enumerate}
\item $\Gamma$ सुसंगत असतो, आणि
\item जर $\Gamma \subsetneq \Gamma'$ असेल, तर $\Gamma'$ विसंगत असतो.
\end{enumerate}
\end{defn}

\begin{explain}
वरील व्याख्येशी समतुल्य अशी पर्यायी व्याख्या ही: वाक्यांचा संच $\Gamma$
\emph{महत्तम सुसंगत} तेव्हा आणि केवळ तेव्हाच असतो, जेव्हा
\begin{enumerate}
\item $\Gamma$ सुसंगत असतो, आणि
\item जर $\Gamma \cup \{ !A \}$ सुसंगत असेल, तर $!A \in \Gamma$ असते.
\end{enumerate}
म्हणजेच~$\Gamma$ मध्ये आधीपासून नसलेली !!{sentence}s महत्तम सुसंगत
संच~$\Gamma$ मध्ये जोडली, तर मिळणारा मोठा संच विसंगत होतो.
\end{explain}

\begin{explain}
संपूर्णतेच्या सिद्धतेत महत्तम सुसंगत संच महत्त्वाचे आहेत. कारण प्रत्येक
सुसंगत !!{sentence}s चा संच~$\Gamma$ हा एखाद्या महत्तम सुसंगत
संचात~$\Gamma^*$ सामावतो, याची खात्री देता येते. शिवाय प्रत्येक
!!{sentence}~$!A$ साठी महत्तम सुसंगत संचात $!A$ किंवा त्याचा
निषेध $ \lnot !A$ यांपैकी एखादे वाक्य असते. विशेषतः अणुरूप !!{sentence}s साठी हे खरे
आहे. म्हणून !!{constant}s घालून योग्य रीतीने विस्तारलेल्या भाषेतील
महत्तम सुसंगत संचापासून आपण !!a{structure} रचू शकतो: कोणती अणुरूप
!!{sentence}s~$\Gamma^*$ मध्ये आहेत यावरून !!{predicate}s चे
अर्थनिर्वचन ठरवतो. मग या !!{structure} मध्ये $\Gamma^*$ मधील
सर्व !!{sentence}s (आणि त्यामुळे $\Gamma$ मधीलसुद्धा) सत्य ठरतात,
हे दाखवता येते. या शेवटच्या विधानाच्या सिद्धतेसाठी
$\lnot !A \in \Gamma^*$ तेव्हा आणि केवळ तेव्हाच $!A \notin
\Gamma^*$, तसेच $(!A \lor !B) \in \Gamma^*$ तेव्हा आणि केवळ तेव्हाच
$!A \in \Gamma^*$ किंवा $!B \in \Gamma^*$, इत्यादी गुणधर्म लागतात.
\end{explain}

\begin{prop}
\ollabel{prop:mcs}
$\Gamma$ महत्तम सुसंगत आहे असे माना. मग:
\begin{enumerate}
\item \ollabel{prop:mcs-prov-in} जर $\Gamma \Proves !A$ असेल, तर $!A \in
  \Gamma$.

\item \ollabel{prop:mcs-either-or} कोणत्याही $!A$ साठी $!A \in
  \Gamma$ किंवा $\lnot !A \in \Gamma$.

\tagitem{prvAnd}{\ollabel{prop:mcs-and} $(!A \land !B) \in \Gamma$
  तेव्हा आणि केवळ तेव्हाच $!A \in \Gamma$ आणि $!B \in \Gamma$ ही दोन्ही.}{}

\tagitem{prvOr}{\ollabel{prop:mcs-or} $(!A \lor !B) \in \Gamma$
  तेव्हा आणि केवळ तेव्हाच $!A \in \Gamma$ किंवा $!B \in \Gamma$.}{}

\tagitem{prvIf}{\ollabel{prop:mcs-if} $(!A \lif !B) \in \Gamma$
  तेव्हा आणि केवळ तेव्हाच $!A \notin \Gamma$ किंवा $!B \in \Gamma$.}{}
\end{enumerate}
\end{prop}

\begin{proof}
पुढील सर्व भागांत $\Gamma$ महत्तम सुसंगत आहे असे मानू.
\begin{enumerate}
\item जर $\Gamma \Proves !A$ असेल, तर $!A \in \Gamma$.

$\Gamma \Proves !A$ असे माना. उलट $!A \notin \Gamma$ असेही
मानल्यास, $\Gamma$ महत्तम सुसंगत असल्यामुळे $\Gamma \cup \{!A\}$
विसंगत असतो; म्हणजे $\Gamma \cup \{!A\} \Proves \lfalse$.
\tagrefs{prfSC/{fol:seq:prv:prop:provability-contr},prfND/{fol:ntd:prv:prop:provability-contr}}
वरून $\Gamma$ विसंगत ठरतो. पण $\Gamma$ सुसंगत आहे या गृहीतकाला
हा व्याघात आहे. म्हणून $!A \notin \Gamma$ शक्य नाही; म्हणजे
$!A \in \Gamma$.

\item कोणत्याही $!A$ साठी $!A \in \Gamma$ किंवा $\lnot !A \in \Gamma$.

उलट, एखाद्या~$!A$ साठी $!A \notin \Gamma$ आणि
$\lnot !A \notin \Gamma$ असे दोन्ही आहेत असे माना. $\Gamma$ महत्तम
सुसंगत असल्यामुळे $\Gamma \cup \{!A\}$ आणि $\Gamma \cup \{\lnot !A\}$
दोन्ही विसंगत आहेत. म्हणून $\Gamma \cup \{!A\} \Proves \lfalse$
आणि $\Gamma \cup \{\lnot !A\} \Proves \lfalse$.
\tagrefs{prfSC/{fol:seq:prv:prop:provability-exhaustive},prfND/{fol:ntd:prv:prop:provability-exhaustive}}
वरून $\Gamma$ विसंगत ठरतो, हा व्याघात. म्हणून असे !!a{sentence}~$!A$
असू शकत नाही. प्रत्येक $!A$ साठी $!A \in \Gamma$ किंवा
$\lnot !A \in \Gamma$.

\tagitem{defAnd}{}{%
\iftag{probAnd}{सराव.}{%
$(!A \land !B) \in \Gamma$ तेव्हा आणि केवळ तेव्हाच $!A \in \Gamma$
आणि $!B \in \Gamma$ ही दोन्ही:

पुढील दिशेसाठी $(!A \land !B) \in \Gamma$ असे माना. मग
$\Gamma \Proves !A \land !B$. आता
\tagrefs{prfSC/{fol:seq:prv:prop:provability-land-left},prfND/{fol:ntd:prv:prop:provability-land-left}}
वरून $\Gamma \Proves !A$ आणि $\Gamma \Proves !B$. तसेच
\olref{prop:mcs-prov-in} वरून अपेक्षेप्रमाणे $!A \in \Gamma$
आणि $!B \in \Gamma$.

उलट दिशेसाठी $!A \in \Gamma$ आणि $!B \in \Gamma$ असे माना.
मग $\Gamma \Proves !A$ आणि $\Gamma \Proves !B$. आता
\tagrefs{prfSC/{fol:seq:prv:prop:provability-land-right},prfND/{fol:ntd:prv:prop:provability-land-right}}
वरून $\Gamma \Proves !A \land !B$. \olref{prop:mcs-prov-in}
वरून $(!A \land !B) \in \Gamma$.}}

\tagitem{defOr}{}{%
\iftag{probOr}{सराव.}{%
$(!A \lor !B) \in \Gamma$ तेव्हा आणि केवळ तेव्हाच $!A \in \Gamma$
किंवा $!B \in \Gamma$.

पुढील दिशेचा प्रतिव्यत्यास सिद्ध करण्यासाठी $!A \notin \Gamma$
आणि $!B \notin \Gamma$ असे माना. आपल्याला $(!A \lor !B)
\notin \Gamma$ दाखवायचे आहे. $\Gamma$ महत्तम सुसंगत असल्यामुळे
$\Gamma \cup \{!A\} \Proves \lfalse$ आणि
$\Gamma \cup \{!B\} \Proves \lfalse$. आता
\tagrefs{prfSC/{fol:seq:prv:prop:provability-lor-left},prfND/{fol:ntd:prv:prop:provability-lor-left}}
वरून $\Gamma \cup \{(!A \lor !B)\}$ विसंगत ठरतो. म्हणून
अपेक्षेप्रमाणे $(!A \lor !B) \notin \Gamma$.

उलट दिशेसाठी $!A \in \Gamma$ किंवा $!B \in \Gamma$ असे
माना. मग $\Gamma \Proves !A$ किंवा $\Gamma \Proves !B$. आता
\tagrefs{prfSC/{fol:seq:prv:prop:provability-lor-right},prfND/{fol:ntd:prv:prop:provability-lor-right}}
वरून $\Gamma \Proves !A \lor !B$. \olref{prop:mcs-prov-in}
वरून अपेक्षेप्रमाणे $(!A \lor !B) \in \Gamma$.}}

\tagitem{defIf}{}{%
\iftag{probIf}{सराव.}{%
$(!A \lif !B) \in \Gamma$ तेव्हा आणि केवळ तेव्हाच
$!A \notin \Gamma$ किंवा $!B \in \Gamma$:

पुढील दिशेसाठी $(!A \lif !B) \in \Gamma$ असे माना.
उलट $!A \in \Gamma$ आणि $!B \notin \Gamma$ असेही माना.
या गृहीतकांवरून $\Gamma \Proves !A \lif !B$ आणि
$\Gamma \Proves !A$. आता
\tagrefs{prfSC/{fol:seq:prv:prop:provability-mp},prfND/{fol:ntd:prv:prop:provability-mp}}
वरून $\Gamma \Proves !B$. परंतु \olref{prop:mcs-prov-in} वरून
$!B \in \Gamma$, हा $!B \notin \Gamma$ या गृहीतकाला व्याघात.

उलट दिशेसाठी प्रथम $!A \notin \Gamma$ हे प्रकरण पाहू.
\olref{prop:mcs-either-or} वरून $\lnot !A \in \Gamma$ आणि म्हणून
$\Gamma \Proves \lnot !A$. आता
\tagrefs{prfSC/{fol:seq:prv:prop:provability-lif},prfND/{fol:ntd:prv:prop:provability-lif}}
वरून $\Gamma \Proves !A \lif !B$.
पुन्हा \olref{prop:mcs-prov-in} वरून अपेक्षेप्रमाणे
$(!A \lif !B) \in \Gamma$.

आता $!B \in \Gamma$ हे प्रकरण पाहू. मग $\Gamma \Proves !B$,
आणि
\tagrefs{prfSC/{fol:seq:prv:prop:provability-lif},prfND/{fol:ntd:prv:prop:provability-lif}}
वरून $\Gamma \Proves !A \lif !B$. \olref{prop:mcs-prov-in}
वरून $(!A \lif !B) \in \Gamma$.}}
\end{enumerate}
\end{proof}

\begin{prob}
\olref[fol][com][mcs]{prop:mcs} ची सिद्धता पूर्ण करा.
\end{prob}

\end{document}
